\section{Prime extraction}\label{ex:section}

Put $N_u=2u+1$ and $a_x(u)=A(u)\Psi(u/x)$, so that
$a_x(u)\ge0$ and $X_A=\sum_u a_x(u)$. Throughout this section the
band geometry, the even integer $r_0$, the damping $q_0=1/2$, and
$\Psi$ are the fixed data of the weight construction. In particular,
Lemma~\ref{wt:mass} gives
\begin{equation}\label{ex:mass-reminder}
 X_A\asymp\frac{x}{L}H_QH_P,\qquad X_A\gg xL^{-C_A},\qquad H_P\ge1.
\end{equation}
We use the Type II and congruence distribution results proved above,
together with the following explicitly stated sieve and proxy inputs.

\subsection{Sieve and proxy inputs}

\begin{lemma}[Block sieve, imported]\label{ex:block-sieve}
Consider finitely many objects with nonnegative weights and, at some
primes $p\le z$, designated bad conditions. Suppose that for every
squarefree product $d$ of these primes, including $d=1$, the weight
on which all conditions indexed by $p\mid d$ hold is
$Xg(d)+r(d)$. Here $X\ge0$, $g$ is multiplicative, and, for fixed
$\eta_0>0$ and $C$, one has
\[
 0\le g(p)\le1-\eta_0,\qquad
 \sum_{w<p\le w^2}g(p)\le C\quad(w>1).
\]
For every sufficiently large even integer $h$, the weight avoiding
all the bad conditions equals
\begin{equation}\label{ex:block-formula}
 X\prod_{p\le z}(1-g(p))(1+O(e^{-h}))
 +O\!\left(\sum_{\substack{d\le z^{4h+2}\\d\text{ squarefree}}}
                 |r(d)|\right).
\end{equation}
The products and sums use only the designated primes. Constants
depend only on the two density bounds, and are uniform when $h$
grows. For $h=2$ there is an upper bound by a fixed constant times
the displayed main term, plus the same remainder sum.

The upper and lower inclusion--exclusion polynomials have
coefficients of absolute value at most one, supported on squarefree
$d\le z^{4h+2}$. The upper polynomial is nonnegative. Its overcount
is bounded by the nonnegative difference between the upper and
lower polynomials; that difference also has coefficients of absolute
value at most one and the same support bound.
\end{lemma}

This is the Block sieve, Lemma~2.9 of \cite{SmoothShifted}. Its
coefficient bound, its growing-$h$ uniformity, and its fixed-depth
upper bound will each be used below. No sieve assertion about the
present weight is included in the import: its required remainder
estimates come from Theorem~\ref{cg:distribution} or from the explicit
box count below.

Define
\begin{equation}\label{ex:euler-products}
 \begin{split}
 V(y)&=\prod_{p\le y}(1-1/p),\qquad
 V_1(y)=\prod_{3\le p\le y}(1-1/(p-1)),\\
 \mathfrak S&=2\prod_{p\ge3}\left(1-\frac1{(p-1)^2}\right)>0.
 \end{split}
\end{equation}
Mertens' estimate and the convergent quotient product give
\begin{equation}\label{ex:mertens}
 V(y)\sim\frac{e^{-\gamma_E}}{\log y},\qquad
 V_1(y)\sim\mathfrak S V(y),
\end{equation}
where $\gamma_E$ is Euler's constant.

Here are the proxy estimates we use. For the moment let
$0<\kappa<1/50$ and $0<b_1<1/2-\kappa$ be arbitrary fixed numbers,
and let $(\gamma,\gamma']$ range over a fixed finite partition of
$(b_1,1/2-\kappa]$. On
$[x^{b_1/4},x^{1-b_1/4}]$ consider the tests
\[
 I_{\mathrm{pr}}(m)=\1_{m\text{ prime}},\qquad
 I_\gamma(m)=\1_{P^-(m)>x^\gamma}.
\]
For a fixed precision $S_1>0$, partition the logarithmic axis into
cells of width $\Delta_L=L^{-S_1}$. On every full cell
$C_Y=(Y,e^{\Delta_L}Y]$ meeting the ambient interval, put
\begin{equation}\label{ex:proxy-definition}
 c_I(C_Y)=
 \frac{\sum_{m\in C_Y}I(m)}
      {\#\{m\in C_Y:P^-(m)>W\}},\qquad
 B_I(m)=c_I(C_Y)\1_{P^-(m)>W}\quad(m\in C_Y).
\end{equation}
Even if a later factor range cuts a cell, both counts in this
definition are taken on the full cell. Range restrictions are
applied only after the ratio has been defined.

\begin{lemma}[Proxy estimates, imported]\label{ex:proxies}
For sufficiently large $x$, the denominators in
\eqref{ex:proxy-definition} are positive and $0\le c_I\le1$.
For every fixed $B>0$, every fixed $S_1>2B+3$, and every interval
$K$ inside a dyad in the stated factor ranges, the sequence
\[
 \alpha_m=(I(m)-B_I(m))\1_{m\in K}
\]
vanishes outside $P^-(m)>W$, has absolute value at most one, and
satisfies
\begin{equation}\label{ex:proxy-discrepancy}
 \left|\sum_{m\in J}\alpha_m\chi(m)m^{-1+it}\right|\le L^{-B}
\end{equation}
for every interval $J\subset[1,x^2]$, every Dirichlet character
of modulus at most $L^B$, and every $|t|\le L^B$.

Write
\[
 \mathcal B_\kappa=[x^{1/2-2\kappa},x^{1/2+2\kappa}].
\]
There is an absolute constant $C_{\mathrm{pr}}$ such that
\begin{equation}\label{ex:prime-density}
 c_{\mathrm{pr}}(m)\le\frac{C_{\mathrm{pr}}}{LV(W)}
       \qquad(m\in\mathcal B_\kappa).
\end{equation}
The constant is independent of $\kappa$ and $S_1$; the threshold
for $x$ may depend on all fixed parameters.

For $w>\gamma>0$, set
\begin{equation}\label{ex:rough-density}
 D_\gamma(w)=\frac1w+
 \sum_{l\ge2}\frac1{l!}
 \int_{\substack{t_i\ge\gamma\ (1\le i<l)\\
                   \sum_{i<l}t_i\le w-\gamma}}
       \frac{\prod_{i<l}(dt_i/t_i)}{w-\sum_{i<l}t_i}.
\end{equation}
This sum is locally finite and nonnegative, and is jointly
continuous on compact subsets of $w>\gamma>0$, including the
thresholds $w=l\gamma$. If $mn=N_u$ on the support of $\Psi(u/x)$
and $\log n/L\in(\gamma,\gamma']$, then
\begin{equation}\label{ex:rough-proxy-bound}
 c_\gamma(m)\le
 \frac{\sup_{\gamma\le t\le\gamma'}D_\gamma(1-t)+o(1)}{LV(W)}.
\end{equation}
Here the parameters are restricted to a compact subset of
$w>\gamma>0$, as holds for the bins under consideration. For fixed
sufficiently small $b>0$,
\begin{align}
 \int_b^{1/2}D_\alpha(1-\alpha)\frac{d\alpha}{\alpha}
     &=D_b(1)-1,\label{ex:density-integral}\\
 D_b(1)&\le\frac{e^{-\gamma_E}}b(1+O(e^{-c/b}))
       \label{ex:density-upper}
\end{align}
with an absolute $c>0$. The endpoint of the integral is interpreted
by its left limit.
\end{lemma}

The discrepancy assertion and the density assertions are,
respectively, Lemmas~7.3 and~7.4 of \cite{SmoothShifted}, in the
section ``Extracting primes with smooth predecessors.'' These
results concern ordinary integer cells; their hypotheses contain
no candidate weight. In particular they apply to the present $A$.
The exact bound $c_I\le1$ also follows directly because each test
selects a subset of the $W$-rough integers in its cell. Thus the
coefficient bounds of both $B_I$ and $I-B_I$ do not depend on the
precision $S_1$. In \eqref{ex:rough-proxy-bound},
\[
 \frac{\log m}{L}=1-\frac{\log n}{L}+O(1/L),\qquad
 \frac{\log m}{L}-\gamma\ge2\kappa+O(1/L),
\]
which verifies the required separation from $w=\gamma$. Every
little-oh assertion in the proxy estimates is taken only after its
parameters have been fixed.

\subsection{A uniform positive pair bound}

We first prove an upper bound that uses only the original weight
geometry. This establishes the constant needed to choose $\kappa$
without referring to $b_1$, the Type II discrepancy exponent, or
the eventual proxy precision.

Choose a sufficiently large fixed integer $j'$ and a fixed $\theta$
in the gap
\begin{equation}\label{ex:micro-choice}
 c_2s'2^{-j'}<\theta<c_1s'2^{-j'+1},\qquad
 10\theta\le\frac1{20},\qquad
 \sum_{j\ge j'}r_0c_2s'2^{-j}\le\frac1{25}.
\end{equation}
Such choices exist because $c_2<2c_1$, and all endpoints decrease
geometrically. Call $Q_j$ with $j\ge j'$ the micro bands and those
with $j<j'$ the macro bands.

\begin{lemma}[Assignment and mask majorant]\label{ex:mask}
An assignment $a$ consists of the complete ordered $r_0$-tuples in
all the micro bands and one prime mark from each $\mathcal P_g$.
Let $D_1=D_1(a)$ be their product, with multiplicities retained, and
put
\[
 \lambda_a=\prod_{j=j'}^{j_x}\frac1{r_0!}\prod_g V_g^{-1}.
\]
For sufficiently large $x$,
\begin{equation}\label{ex:assignment-masses}
 D_1\le x^{1/20},\qquad
 \sum_a\frac{\lambda_a}{D_1}\ll H_Q,\qquad
 \sum_a\lambda_a\ll x^{1/20}H_Q.
\end{equation}
For each assignment, ban every odd non-$\mathcal P$ prime at most
$x^\theta$ not dividing $D_1$, and independently ban each
$\mathcal P$ prime not dividing $D_1$ with probability $1-q_0$.
With expectation over this finite random mask,
\begin{equation}\label{ex:positive-majorant}
 A(u)\le A_0(u)\le
 \sum_a\lambda_a\1_{D_1\mid u}\,\mathbb E_a
           \1_{u\text{ is divisible by no banned prime}}.
\end{equation}
The constants in \eqref{ex:assignment-masses} depend only on the
fixed weight geometry and $j'$.
\end{lemma}

\begin{proof}
The product of the micro entries is at most $x^{1/25}$ by
\eqref{ex:micro-choice}. There are $O(T)$ marks and every mark is at
most $\exp(L^{2/5})$, so their product is
$\exp(O(TL^{2/5}))=x^{o(1)}$. This proves the asserted bound for
$D_1$, including assignments with repeated micro entries.
Write
\[
 H_Q^{\mathrm{mic}}
 =\prod_{j=j'}^{j_x}\frac1{r_0!}
                      \left(\sum_{p\in Q_j}\frac1p\right)^{r_0}.
\]
Summing the reciprocal product of the marks gives
$\prod_g(V_g^{-1}\sum_{p\in\mathcal P_g}p^{-1})=1$.
Consequently $\sum_a\lambda_a/D_1=H_Q^{\mathrm{mic}}\ll H_Q$:
only finitely many fixed band factors, bounded above and below,
have been omitted from $H_Q$. The last assertion in
\eqref{ex:assignment-masses} follows by multiplying by the bound
for $D_1$.

To prove the majorant, fix $u$ with $A_0(u)>0$ and a compatible
choice of ordered micro entries and marks. The band gap in
\eqref{ex:micro-choice} implies that all non-$\mathcal P$ prime
factors of $u$ below $x^\theta$ are among these micro entries,
whereas every macro prime exceeds $x^\theta$. All group primes
are below $x^\theta$ for sufficiently large $x$. The only random
survival conditions are therefore at the distinct unmarked
group-prime divisors of $u$. There are exactly
$\omega_P(u)-s_P$ of them, so the survival probability is
$q_0^{\omega_P(u)-s_P}$, even when marked or unmarked primes
occur to higher powers.

Summing one mark per group with normalization $\prod_g V_g^{-1}$
recovers the factor $W_1(u)$. For a fixed prime multiset in any
band, the number of ordered $r_0$-tuples divided by $r_0!$ is
$1/\prod_p v_p!\le1$. Summing the normalized compatible micro
lists retains their exact contribution to $a(u_*)$; dropping the
remaining macro-list contribution can only increase it. Thus
the compatible assignments already dominate $A_0(u)$.
Incompatible assignments contribute nonnegative terms, proving
\eqref{ex:positive-majorant}.
\end{proof}

\begin{lemma}[Positive dyadic pair estimate]\label{ex:box}
There is a constant $C_{\mathrm{box}}$, depending only on the
fixed weight data and the choices in \eqref{ex:micro-choice}, such
that, on every dyadic box with $M_1M_2\asymp x$ and
$M_i\ll x^{0.54}$,
\begin{equation}\label{ex:box-bound}
 \sum_{\substack{M_1\le m<2M_1,\ M_2\le n<2M_2\\
                 mn=N_u,\ P^-(m),P^-(n)>W}}
       a_x(u)\le C_{\mathrm{box}}X_AV(W)^2.
\end{equation}
The fixed constants in the size comparisons can be taken to cover
all boxes meeting $mn=2u+1$ with $u\in[x,2x]$.
\end{lemma}

\begin{proof}
Use Lemma~\ref{ex:mask} and bound $\Psi$ by its fixed supremum.
Fix an assignment and a mask. Since $D_1$ is odd, $D_1\mid u$
implies $mn\equiv1\pmod{D_1}$. We count all pairs in the box with
this congruence, dropping both the parity of $mn$ and the remaining
support restriction on $u$.

For every odd prime $l\le x^\theta$ not dividing $D_1$, impose
the bad condition $mn\equiv0\pmod l$ when $l\le W$, and the bad
condition $mn\equiv1\pmod l$ when $l$ is banned. There are
$2l-1$ pairs of residues with product zero and $l-1$ with product
one; the two sets are disjoint. Hence the local number and density
of bad residue pairs are
\begin{equation}\label{ex:local-density}
 \nu(l)=(2l-1)\1_{l\le W}+(l-1)\1_{l\text{ banned}},
 \qquad g(l)=\frac{\nu(l)}{l^2}.
\end{equation}
In particular $g(l)\le7/9$ and $g(l)\le3/l$. These give uniform
density hypotheses for Lemma~\ref{ex:block-sieve}, independent of
the assignment and mask.

There are exactly $\varphi(D_1)$ residue pairs with product one
modulo $D_1$: choose the first unit and then its unique inverse.
This statement does not require $D_1$ to be squarefree. For a
squarefree $d$ built from the sieving primes, the Chinese remainder
theorem gives exactly $\varphi(D_1)\nu(d)$ pairs modulo $D_1d$,
where $\nu$ is multiplicative. Each pair of classes contributes
\[
 \frac{M_1M_2}{(D_1d)^2}
  +O\left(1+\frac{M_1+M_2}{D_1d}\right).
\]
Thus the base mass is
$M_1M_2\varphi(D_1)/D_1^2$, with local density $g(d)$.
Since $\nu(d)\le d\tau(d)^2$, the lattice remainders satisfy
\[
 |r_a(d)|\ll\tau(d)^2\{M_1+M_2+D_1d\}.
\]
Use the $h=2$ upper bound in Lemma~\ref{ex:block-sieve}, whose
level is $D=(x^\theta)^{10}$. Fixed divisor-moment bounds give
\begin{equation}\label{ex:box-error}
 \sum_{d\le D}|r_a(d)|
 \ll L^{C_0}\{(M_1+M_2)D+D_1D^2\}
\end{equation}
for a fixed $C_0$. This includes $d=1$ and is uniform in the mask.

We next compare the sieve product with Mertens products. For
$l\le W$ banned by the mask, its factor is
$1-3/l+2/l^2$; for an allowed $l\le W$ it is $(1-1/l)^2$.
For $W<l\le x^\theta$ outside $D_1$, all primes are non-group
primes and banned, and the factor is $1-1/l+1/l^2$.
Comparing these factors, allowing a convergent product of
$1+O(l^{-2})$ and the fixed factor at $2$, gives
\begin{equation}\label{ex:product-comparison}
 \prod_{\substack{3\le l\le x^\theta\\l\nmid D_1}}(1-g(l))
 \ll V(W)^2V(x^\theta)
       \prod_{l\mid D_1}(1-1/l)^{-3}
       \prod_{\substack{l\in\mathcal P,\ l\nmid D_1\\
                         l\text{ allowed}}}(1-1/l)^{-1}.
\end{equation}
Every prime dividing $D_1$ is at least $L^{20}$, and
$\log D_1=O(L)$, so
\[
 \sum_{l\mid D_1}\frac1l
       \ll\frac{L^{-19}}{\log L}.
\]
Its correction product in \eqref{ex:product-comparison} is
therefore uniformly bounded. Independence of the mask gives
\begin{align*}
 \mathbb E_a
 \prod_{\substack{l\in\mathcal P,\ l\nmid D_1\\l\text{ allowed}}}
              (1-1/l)^{-1}
 &=\prod_{\substack{l\in\mathcal P\\l\nmid D_1}}
              \left(1+\frac{q_0}{l-1}\right)\\
 &\le\prod_{l\in\mathcal P}\left(1+\frac{q_0}{l-1}\right)
 \le H_P.
\end{align*}
For the last inequality, the group factor in $H_P$ is
\[
 \prod_{p\in\mathcal P_g}\left(1+\frac{q_0}{p-1}\right)
       \frac1{V_g}\sum_{p\in\mathcal P_g}\frac1{p-1+q_0},
\]
and the final factor is at least one.

Since $\varphi(D_1)/D_1^2\le1/D_1$, summing the main term over
assignments with Lemma~\ref{ex:mask} gives
\begin{equation}\label{ex:box-main}
 \ll xV(W)^2V(x^\theta)H_QH_P
 \ll X_AV(W)^2.
\end{equation}
Here $V(x^\theta)\asymp L^{-1}$, with $\theta$ already fixed,
and we used \eqref{ex:mass-reminder}.
It remains essential to sum the errors over assignments as well.
From \eqref{ex:assignment-masses},
\[
 \sum_a\lambda_aD_1\le x^{1/10}O(H_Q).
\]
Since $D\le x^{1/20}$ and $M_i\ll x^{0.54}$,
\begin{align}
 \sum_a\lambda_a\sum_{d\le D}|r_a(d)|
 &\ll L^{C_0}H_Q
       \{x^{1/20}(M_1+M_2)D+x^{1/10}D^2\}\notag\\
 &\ll L^{C_0}H_Q(x^{0.64}+x^{0.20})
   =o(X_AV(W)^2).\label{ex:all-box-errors}
\end{align}
Indeed $X_AV(W)^2\asymp xH_QH_P/L^2$ and $H_P\ge1$.
This proves \eqref{ex:box-bound}. Every constant used here depends
only on the original weight data and $j',\theta$.
\end{proof}

For later use, a dyadic decomposition of $\mathcal B_\kappa$ and
Lemma~\ref{ex:box} imply
\begin{equation}\label{ex:all-balanced-pairs}
 R_\kappa:=
 \sum_{\substack{mn=N_u,\ m,n\in\mathcal B_\kappa\\
                  P^-(m),P^-(n)>W}}a_x(u)
 \le C_1(\kappa L+1)X_AV(W)^2.
\end{equation}
There are $O(\kappa L+1)$ first-factor dyads because the interval
has logarithmic length $4\kappa L$, and only $O(1)$ second-factor
dyads per first-factor dyad because $mn\in[2x+1,4x+1]$.
Enlarging partially intersected boxes is legitimate by positivity.
For $\kappa<1/50$ all their sides are $O(x^{0.54})$.
Consequently $C_1$ is independent of $\kappa,b_1,S_1$.
Fix, at this point, the constant
\begin{equation}\label{ex:balanced-constant}
 C_{\mathrm{bal}}=C_{\mathrm{pr}}^2C_1,
\end{equation}
enlarging it if necessary to absorb the fixed dyadic conventions.

\subsection{Parameter choice and test replacement}

Choose $0<\kappa<1/50$ so small that
\begin{equation}\label{ex:kappa-choice}
 C_{\mathrm{bal}}\kappa<\mathfrak S/4.
\end{equation}
Next choose $b_1>0$ sufficiently small, with $b_1<1/100$, to
satisfy the presieve and subtraction conditions below. The
constants in those conditions do not depend on a later proxy
precision. Set
\[
 \vartheta=\frac12-\frac\kappa2,\qquad b_* = b_1/3.
\]
By continuity in Lemma~\ref{ex:proxies}, a sufficiently fine fixed
partition of $(b_1,1/2-\kappa]$ satisfies
\begin{equation}\label{ex:mesh}
 \sum_{(\gamma,\gamma']}
    \sup_{\gamma\le t\le\gamma'}D_\gamma(1-t)
                      \log(\gamma'/\gamma)
 \le D_{b_1}(1)-1+\frac1{10}.
\end{equation}
In fact these sums converge to
$\int_{b_1}^{1/2-\kappa}D_\alpha(1-\alpha)\,d\alpha/\alpha$:
on this fixed compact range, replacing $\gamma$ and $t$ by the
same point changes the continuous integrand uniformly by a
quantity tending to zero with the mesh. Positivity and
\eqref{ex:density-integral} then give \eqref{ex:mesh}.

Prescribe a Type II saving $D_*>C_A+4$ and a congruence saving
$D_{\mathrm{cg}}>C_A+4$. Apply Theorem~\ref{ti:distribution} with
$b_*$ and coefficient bound one, obtaining its required fixed
discrepancy exponent $B$. Only now fix $S_1>2B+3$ and define the
proxies by \eqref{ex:proxy-definition}. All these parameters,
including any prime number theorem accuracies in the proxy inputs,
are fixed before $x$ tends to infinity.

\begin{lemma}[Replacement on the required factor ranges]
\label{ex:replacement}
In all factor sums used below, a test $I$ on either factor of
$mn=N_u$ can be replaced by $B_I$ with total error $o(X_A/L)$,
provided the coefficient on the other factor has absolute value at
most one. This holds also when that other coefficient is a proxy.
\end{lemma}

\begin{proof}
For the unbalanced sums, the prime factor $n$ is in
$(x^{b_1},x^{1/2-\kappa}]$, and the complementary factor lies
between fixed positive multiples of $x^{1/2+\kappa}$ and
$x^{1-b_1}$. For the balanced sums both factors are in
$\mathcal B_\kappa$. Thus every contributing factor dyad, after
intersecting with its actual interval, has scale between
$x^{b_1/3}$ and $x^{1-b_1/3}$ for sufficiently large $x$.
The actual factor ranges also lie inside the ambient proxy
interval. The fixed gaps in these exponent inequalities absorb
all dyadic endpoint constants.

On a dyad use $(I-B_I)\1_K$ as the first sequence in
Theorem~\ref{ti:distribution}, where $K$ is the actual summation
interval. Lemma~\ref{ex:proxies} supplies its rough support,
coefficient bound, and discrepancy. Exchange the factor names
when the test is on the second factor. The theorem requires only
boundedness of the companion sequence, so a nonsmooth proxy is
permitted. There are $O(L)$ contributing dyadic pairs for each
of the fixed finitely many tests. The total error is therefore
$O(xL^{-D_*+1})=o(X_A/L)$ by \eqref{ex:mass-reminder}.
The full-cell normalization is retained when $K$ cuts a cell.
In particular, neither $B$ nor the coefficient exponent used in
the Type II theorem depends on $S_1$.
\end{proof}

\subsection{Presieving and unbalanced composites}

For odd squarefree $d$ write, as in the congruence theorem,
\[
 r(d)=\sum_{d\mid N_u}a_x(u)-\frac{X_A}{\varphi(d)}.
\]
Theorem~\ref{cg:distribution}, with the already fixed
$\vartheta,D_{\mathrm{cg}}$, gives
\begin{equation}\label{ex:remainder-sum}
 \sum_{\substack{d\le x^\vartheta\\d\text{ odd and squarefree}}}
       |r(d)|
 \ll xL^{-D_{\mathrm{cg}}}+X_AL^{-18}=o(X_A/L).
\end{equation}

\begin{lemma}[Initial sifted mass]\label{ex:presieve}
For sufficiently small fixed $b_1$,
\begin{equation}\label{ex:presieve-bound}
 \sum_{P^-(N_u)>x^{b_1}}a_x(u)
 \ge\frac{\mathfrak S X_A}{L}
       \left\{\frac{e^{-\gamma_E}}{b_1}
                    (1-O(e^{-c/b_1}))+o(1)\right\}.
\end{equation}
\end{lemma}

\begin{proof}
Apply Lemma~\ref{ex:block-sieve} to the nonnegative weights
$a_x(u)$, with bad condition $p\mid N_u$ at odd primes
$p\le z=x^{b_1}$. Its density is $g(p)=1/(p-1)$, at most $1/2$;
the sums of these densities over $(w,w^2]$ are bounded. Take
\[
 h=2\left\lfloor\frac1{40b_1}\right\rfloor.
\]
Making $b_1$ small ensures that this even integer is sufficiently
large for the two-sided sieve. Its level satisfies
\[
 b_1(4h+2)\le\frac15+2b_1<\vartheta.
\]
Thus \eqref{ex:remainder-sum} supplies the full remainder sum.
Since $N_u$ is odd, avoiding the specified odd primes is exactly
$P^-(N_u)>x^{b_1}$. Finally use \eqref{ex:mertens} and
$e^{-h}=O(e^{-c/b_1})$ in \eqref{ex:block-formula}.
\end{proof}

\begin{lemma}[Conditional small-prime sieve]\label{ex:conditional}
For each fixed bin $(\gamma,\gamma']$ of the chosen partition,
\begin{equation}\label{ex:conditional-sieve}
 \sum_{\substack{x^\gamma<n\le x^{\gamma'}\\n\text{ prime}}}
       \sum_{\substack{n\mid N_u\\P^-(N_u)>W}}a_x(u)
 =X_AV_1(W)\{\log(\gamma'/\gamma)+o(1)\}.
\end{equation}
\end{lemma}

\begin{proof}
Fix such a prime $n$. For each odd squarefree $d$ supported on
primes at most $W$, one has $n>W$ for sufficiently large $x$, and
\[
 \sum_{nd\mid N_u}a_x(u)
   =\frac{X_A}{(n-1)\varphi(d)}+r(nd).
\]
Apply Lemma~\ref{ex:block-sieve} on the objects with $n\mid N_u$,
using base mass $X_A/(n-1)$, local density $1/(p-1)$, and
$h=2\lceil T^2\rceil$. Its stated growing-$h$ uniformity applies;
the level for $d$ is
\[
 D_W=W^{4h+2}=\exp(O(\sqrt L\,T^2))=x^{o(1)}.
\]
Every product $nd$ used is at most
$x^{1/2-\kappa+o(1)}<x^\vartheta$.

Crucially, the map $(n,d)\mapsto nd$ is injective, since $n$ is
the unique prime divisor of $nd$ exceeding $W$. This remains true
when the sums are taken over all the disjoint bins. The sieve
polynomials have coefficients of absolute value at most one, so
the sum of the remainders over all $n$ and $d$ is bounded by the
single sum \eqref{ex:remainder-sum}, with no divisor multiplicity.
It is $o(X_AV_1(W))$, since $V_1(W)\asymp L^{-1/2}$.
The main terms are
\[
 X_AV_1(W)(1+O(e^{-h}))
          \sum_{\substack{x^\gamma<n\le x^{\gamma'}\\n\text{ prime}}}
                         \frac1{n-1}.
\]
Partial summation from the prime number theorem gives the last
sum as $\log(\gamma'/\gamma)+o(1)$, proving the result.
\end{proof}

\begin{lemma}[Removal of unbalanced composites]\label{ex:unbalanced}
The weight of prime values $N_u$, together with composite values
whose least prime factor exceeds $x^{1/2-\kappa}$, is at least
\begin{equation}\label{ex:surviving-mass}
 \frac{\mathfrak S X_A}{2L}
\end{equation}
for sufficiently large $x$ after $b_1$ is chosen sufficiently small.
\end{lemma}

\begin{proof}
If a presieved composite has least prime factor
$n\in(x^\gamma,x^{\gamma'}]$, every prime factor of $m=N_u/n$
is at least $n>x^\gamma$. This includes a repeated occurrence of
$n$. Its weight is therefore bounded above by the sum over
$mn=N_u$ of $I_\gamma(m)\1_{x^\gamma<n\le x^{\gamma'},\ n\text{ prime}}
a_x(u)$. Other divisors counted by this sum merely increase the
upper bound.

Replace $I_\gamma$ by $B_\gamma$ using
Lemma~\ref{ex:replacement}, and use \eqref{ex:rough-proxy-bound}.
Because $n>W$, $P^-(m)>W$ is equivalent to $P^-(N_u)>W$; no
coprimality between $m$ and $n$ is needed. Lemma~\ref{ex:conditional}
and \eqref{ex:mertens} bound the weight removed for this bin by
\[
 \frac{\mathfrak S X_A}{L}
 \left\{\sup_{\gamma\le t\le\gamma'}D_\gamma(1-t)
                       \log(\gamma'/\gamma)+o(1)\right\},
\]
apart from the total $o(X_A/L)$ replacement error. Summing over
the fixed mesh and using \eqref{ex:mesh} bounds the total removed
weight by
\[
 \frac{\mathfrak S X_A}{L}
       \{D_{b_1}(1)-1+1/10+o(1)\}.
\]
Subtract this from Lemma~\ref{ex:presieve}. By
\eqref{ex:density-upper}, the coefficient of
$\mathfrak S X_A/L$ left over is at least
\[
 1-\frac1{10}-O(b_1^{-1}e^{-c/b_1})+o(1).
\]
Choose $b_1$ so small that this is greater than $1/2$ for
sufficiently large $x$. This choice also satisfies the initial
sieve conditions. The surviving presieved values are exactly
among the two classes in the assertion, so their weight proves
\eqref{ex:surviving-mass}.
\end{proof}

\subsection{Balanced composites and the conclusion}

\begin{lemma}[Balanced composite bound]\label{ex:balanced}
With the constant fixed in \eqref{ex:balanced-constant}, the weight
of composite $N_u$ with $P^-(N_u)>x^{1/2-\kappa}$ is at most
\begin{equation}\label{ex:balanced-bound}
 C_{\mathrm{bal}}(\kappa+1/L)\frac{X_A}{L}+o(X_A/L).
\end{equation}
\end{lemma}

\begin{proof}
Such a composite has exactly two prime factors counted with
multiplicity: three would give
$N_u>x^{3(1/2-\kappa)}>4x+1$. Both factors lie in
$\mathcal B_\kappa$ for sufficiently large $x$, since the larger
one is at most $(4x+1)x^{-1/2+\kappa}<x^{1/2+2\kappa}$.
Ordered prime pairs cover every such composite at least once;
distinct primes give two pairs and a repeated prime gives one.
Thus their total weight is bounded by
\[
 \sum_{\substack{mn=N_u\\m,n\in\mathcal B_\kappa}}
              I_{\mathrm{pr}}(m)I_{\mathrm{pr}}(n)a_x(u).
\]
Use the exact identity
\[
 I_{\mathrm{pr}}(m)I_{\mathrm{pr}}(n)
       -B_{\mathrm{pr}}(m)B_{\mathrm{pr}}(n)
 =(I_{\mathrm{pr}}-B_{\mathrm{pr}})(m)I_{\mathrm{pr}}(n)
   +B_{\mathrm{pr}}(m)(I_{\mathrm{pr}}-B_{\mathrm{pr}})(n).
\]
Lemma~\ref{ex:replacement} applies to the two terms in turn,
exchanging factor roles in the second. Both companion sequences
have absolute value at most one, independently of $S_1$.
The resulting positive proxy sum is, by
\eqref{ex:prime-density}, at most
\begin{equation}\label{ex:prime-to-positive}
 \frac{C_{\mathrm{pr}}^2}{L^2V(W)^2}
 \sum_{\substack{mn=N_u,\ m,n\in\mathcal B_\kappa\\
                  P^-(m),P^-(n)>W}}a_x(u).
\end{equation}
Apply \eqref{ex:all-balanced-pairs} and
\eqref{ex:balanced-constant}, and add the $o(X_A/L)$ replacement
error. This proves \eqref{ex:balanced-bound}.
\end{proof}

\begin{proof}[Proof of Theorem~\ref{thm:main}]
Combining Lemmas~\ref{ex:unbalanced} and~\ref{ex:balanced} gives
\[
 \sum_{N_u\text{ prime}}a_x(u)
 \ge\left(\frac{\mathfrak S}{2}
               -C_{\mathrm{bal}}\kappa-o(1)\right)\frac{X_A}{L}
 \gg\frac{X_A}{L},
\]
by \eqref{ex:kappa-choice}. Lemma~\ref{wt:squarefree} removes only
$o(X_A/L)$ from this mass.
Consequently there is a supported $u$ with $2u+1$ prime and
$2u$ squarefree.

For completeness, the support condition counts prime factors with
their multiplicities. The ordinary part $u_*$ has
$r_0(j_x+1)$ prime factors, an even number because $r_0$ is even.
The factor $(1-\mathcal E(u))/2$ in $A(u)$ imposes an odd total
multiplicity on the group-prime part. Hence $\Omega(u)$ is odd.
Every supported $u$ is odd, so
$\Omega(2u)=1+\Omega(u)$ is even. On the retained squarefree
predecessors this is also an even number of distinct prime
factors. Finally $u\in[x,2x]$ wherever $a_x(u)>0$; the resulting
primes exceed $x$ for arbitrarily large $x$, proving infinitude.
\end{proof}

The order of choices is intrinsic to the argument: the original
weight data and the micro-band choices determine
$C_{\mathrm{box}}$ and $C_{\mathrm{bal}}$; then come $\kappa$,
$b_1$, the finite mesh, the distribution savings and their
discrepancy exponent, and finally $S_1$. All thresholds for $x$
are imposed afterwards. In particular the absolute prime-proxy
constant and the coefficient bound one, rather than any
precision-dependent regularity, are what make this order valid.
